% FILE: ~/GPLC/MEMORIA/sections/implementacion.tex
\newpage
\section{Implementación}
\label{sec:implementacion}
El enfoque de este capítulo es realizar una presentación que aborda la implementación de las principales estructuras que dan soporte a las operaciones del intérprete, y las rutinas esenciales que operan sobre ellos. Entre las estructuras de datos, se cubren los tokenes de código fuente, ambos AST (fuente y destino), los tipos graduales, los tipos de fórmula y la infraestructura simbólica como parte integral de ellos, y los valores que el sistema retorna. Las principales rutinas tratadas consisten en el parsing de código fuente, operaciones de tipo, resolución de fórmulas, chequeo estático de ambos AST, elaboración,  y evaluación.

\subsection{Estructuras de Datos}
\subsubsection{Tipos Graduales}

La probabilidad gradual en el lenguaje fuente es un componente de la expresión de selección y de los tipos de distribución. No es un tipo en si mismo, pero por se parte integral de los tipos graduales, se implementa en el módulo \verb|Source.Type| como el tipo de dato  algebraico \verb|gprob|, que contempla constructores para la probabilidad concreta y la desconocida. En la sintaxis de OCaml
, el constructor de la probabilidad concreta \verb|Preal| toma un número de punto flotante (la probabilidad) y un entero (un id). El constructor de la probabilidad desconocida \verb|Pany| toma solo un entero, que también representa un id, y es de particular importancia para convertir una probabilidad \verb|?| a una única probabilidad simbólica en distintas instancias de transformación.

Los tipos graduales simples y de distribución están implementados como dos tipos de datos algebraicos mutuamente inductivos: el constructor funcional del tipo simple, \verb|Fun|, toma como argumento el par \verb|(simple, distr)|, y el constructor básico de distr, \verb|Distr|, toma como argumento una lista de pares \verb|(simple, gprob)|.  El resto de los constructores del tipo simple son
\verb|Real|, \verb|Bool| y \verb|Any|, este último representando al tipo desconocido \verb|?|. Los demás constructores del tipo de distribución son \verb|Scale| y
\verb|Sum|, los cuales escalan una distribución por una probabilidad gradual, y suman una lista de distribuciones, respectivamente.

\begin{lstlisting}[language=OCaml,caption={Probabilidad Gradual y Tipos Graduales},label={asdf}]
type gprob  = (* Gradual Probability *)
| Preal of float * int (* Use: Preal(probabilidad: float, id: int)*)
| Pany of int 

type simple = (* Gradual Simple Types *)
| Real | Bool
| Any  | Fun of simple * distr

and distr = (* Gradual Distribution Type *)
| Distr of (simple * gprob) list (* Use: Distr [(t1,p1); ... ; (tn,pn)] *) 
| Scale of gprob * distr 
| Sum   of distr list\end{lstlisting}


\subsubsection{AST fuente}

El arbol de sintaxis abstracta del lenguaje fuente está implementado en el
módulo \verb|Source.Term| como un tipo de datos algebraico inductivo. Este es un
tipo parametrizado por dos variables de tipo \verb|'st| y \verb|'dt|. Sus
constructores modelan la sintaxis abstracta del lenguaje fuente. Cada
constructor a su vez recibe como último argumento un valor de otro tipo
paramétrico: la anotación \verb|'t annot|, en la cual se inserta el tipo del
nodo, una vez inferido por el chequeo estático.

El fundamento de parametrizar el tipo ast con dos variables de tipo, es lograr utilizar
esta misma estructura para modelar los arboles no tipados y tipados: Si se parametriza
ast con el par de tipos (\verb|(unit,unit)|), entonces la anotación hereda el tipo \verb|unit|
y el espeacio reservado para el tipo del nodo se llena con el único valor de este tipo, es decir \verb|()|.
En cambio si el tipo \verb|ast| se parametriza con el par \verb|(simple,distr)|, cada anotación recibirá
uno de estos dos tipos. Cuál de los dos tipos asignar a una anotación determinada es una decisión
que está codificada en la definición del arbol. Por ejemplo, la anotación de la expresión \verb|Num|
es del tipo \verb|'st annot|, lo cual significa que recibirá un tipo simple. En cambio,
la anotación de la expresión \verb|Choice| tiene tipo \verb|'dt annot|, indicando que recibirá un
 tipo de distribución. Este diseño apunta a resolver la tensión entre un arbol no tipado y uno tipado,
y entre la asignación de distintas categorías de tipo a nodos de un mismo arbol.

\begin{lstlisting}[language=OCaml,caption={Arbol de sintaxis abstracta, Lenguaje Fuente},label={lst:lst:src-ast}]
type ('st, 'dt) ast =
| Num       of float * 'st annot
| Bool      of bool * 'st annot 
| Var       of string * 'st annot
| Fun       of string * Type.simple * ('st,'dt) ast * 'st annot
| Sascr     of ('st,'dt) ast * Type.simple * 'st annot
| Dascr     of ('st,'dt) ast * Type.distr * 'dt annot
| Let       of string * ('st,'dt) ast * ('st,'dt) ast list * 'dt annot 
| If        of ('st,'dt) ast * ('st,'dt) ast * ('st,'dt) ast * 'dt annot
| App       of ('st,'dt) ast * ('st,'dt) ast * 'dt annot 
| Choice    of Type.gprob * ('st,'dt) ast * ('st,'dt) ast * 'dt annot 
| NumBinOp  of ('st,'dt) ast * ('st,'dt) ast *  Bin.num_op * 'st annot 
| BoolBinOp of ('st,'dt) ast * ('st,'dt) ast * Bin.bool_op * 'st annot 

type untyped_term = (unit,unit) ast
type typed_term   = (simple,distr) ast\end{lstlisting}

\subsubsection{Fórmulas}
Las fórmulas son sistemas de ecuaciones e inecuaciones, posiblemente no
lineales, cuyas variables corresponden a probabilidades simbólicas de
distribuciones de fórmula.

En la base del concepto de fórmula, está la
probabilidad simbólica, implementada como un constructor \verb|SymProb| que toma
4 argumentos: un primer argumento string, que es la variable de la probabilidad
simbólica propiamente tal, un segundo argumento de tipo entero, que indica la
posición de la probabilidad dentro de la distribución, y finalmente dos
argumentos más de tipo entero, que almacenan, según sea el origen de la
distribución, información acerca de la posicion dentro de la distribución o
información de las posiciones de los tipos simples que dieron origen a la
probabilidad en la distribución producto.

El tipo de dato algebraico que modela las expresiones que se comparan en las
fórmulas se implementa en \verb|Formula.Expr|. Los constructores \verb|Real| y
\verb|Symb| elaboran los casos bases de esta álgebra. El primero toma un número real, y el segundo una probabilidad
simbólica. En particular, el
constructor \verb|Symb| transforma una probabilidad simbólica en una expresión
simbólica. \verb|Add|, \verb|Sub|, \verb|Mul| y \verb|Sub| son las operaciones
reales binarias básicas. \verb|Sum| y \verb|Prod| suman y multiplican listas de
expresiones, respectivamente.

El tipo de datos \verb|Formula.Form.t| es el que representa a las fórmulas como
restricciones, y como conjunciones entre estas mismas.  Los constructores
\verb|Integrity|, \verb|Distribution| y \verb|Null| aplican a una expresión,
restricciones importantes, y por eso reciben nombres significativos. Los
constructores \verb|Eq|, \verb|Lt| y \verb|Le| comparan dos expresiones de con
las relaciones $=$, $<$ y $\leq$. Finalmente, el operador de conjunción
\verb|And| le da su caracter inductivo al tipo de fórmula, representando una
conjunción sobre una lista de fórmulas.



\begin{ocamlblock}[caption={Fórmulas y sus componentes},label={lst:formula}]
  (* Symbolic Probability (module: Sprob) *)
  type t = SymProb of string * int * int * int (* Uso: TagVar("X",5,1,4) *)

  (* Expression (module: Expr) *)
  type t =
  | Real  of float  | Symb  of Tvar.t
  | Add   of t * t  | Sub   of t * t
  | Mul   of t * t  | Div   of t * t
  | Sum   of t list | Prod  of t list

  (* Formula (module: Form) *)
  type t = 
  | Integrity    of Expr.t (* Expression in [0,1] *)
  | Distribution of Expr.t (* Expression = 1.0 *)
  | Null         of Expr.t (* Expression = 0.0*)
  | Eq           of Expr.t * Expr.t (* Expr1 = Expr2 *)
  | Lt           of Expr.t * Expr.t (* Expr1 < Expr2 *)
  | Le           of Expr.t * Expr.t (* Expr1 <= Expr2 *)
  | And          of t list (* And[f1;...;fn] *)
\end{ocamlblock}





\subsubsection{Distribución Paramétrica y Tipos de Fórmula}
En el módulo \verb|Formula.Distribution| se implementa el tipo paramétrico de
distribución \verb|dis|. Este tipo materializa la idea de la distribución de fórmula como
un tipo polimórfico consistente en un par \emph{(distribución, formula)}, donde la 
distribución es una lista de pares de tipo \verb|(t', Sprob.t)|, es decir valores de
tipo \verb|t'| asociados a una probabilidad simólica. La fórmula asociada al
soporte de la distribución impone restricciones a los valores posibles de las
probabilidades de tipo \verb|Sprob.t|. El tipo \verb|'t| representa el
\emph{dominio} de la distribución.

La implementación de los tipos de fórmula es similar a la de los tipos
graduales. Los tipos simples y de distribución son también definidos de manera
mutumente inductiva. Los tipos simples, producidos por los constructores
\verb|Real|, \verb|Bool|, \verb|Any| y \verb|Fun(simple,distr)| son análogos a
aquellos en los tipos graduales. La diferencia fundamental radica en los tipos de
distribución: estos se definen como el tipo \verb|dis|, parametrizado por el
tipo \verb|simple|. De esta manera, el tipo \verb|distr| es efectivamente
definido como una distribución simbólica de tipos simples. 

\begin{ocamlblock}[caption={Distribución Paramétrica y Tipos de Fórmula}, label={lst:ftype}]
  (* Polymorphic Distribution Type (module: Distribution) *)
  type 't dis = (('t * Sprob.t) list) * Form.t

  (* Simple Formula Types *)
  type simple = 
  | Real
  | Bool
  | Any 
  | Fun of simple * distr

  (* Distribution Formula Types *)
  and distr = simple Distribution.dis
\end{ocamlblock}



\subsubsection{Coupling}

\TODO{}



\subsubsection{AST destino}
El AST destino, implementado en \verb|Target.Term| es muy semejante al AST del
lenguaje fuente, dado que también se define de manera paramétrica para
adaptarse a los requerimientos previos y posteriores al chequeo estático del
lenguaje destino. El componente de tipo \verb|'t annot| es el mismo tipo usado
en los nodos del arbol fuente, dado que el parámetro \verb|'t| puede ser
cualquier tipo (simple o distribución, gradual o de fórmula). Esto facilita
además la propagación de información posicional desde el lenguaje fuente al
destino, ya que los metadatos dentro de la anotación son trasladados por defecto
a su nodo destino durante la elaboración.

En el código \ref{lst:trgt-ast} se muestra un extracto de la
definición de los nodos del arbol destino, incluyendo solo aquellos
que se diferencian de sus nodos correspondientes en el lenguaje
fuente.  Los constructores \verb|Sascr| y \verb|Dascr|, que producen
las expresiones de ascripción simple y de distribución,
respectivamente, ahora reciben un primer argumento adicional, que
constituye la evidencia de la ascripción. El constructor de
\verb|Choice|, a diferencia de su versión fuente, recibe dos
probabilidades simbólicas: la de selección de la primera rama, y su
complemento. El tipo \verb|fprob| es en realidad un tipo definido
como \verb|type fprob = Sprob.t * Form.t|, es decir una probabilidad
simbólica $x$ asociada a una fórmula $x=r$ o $0\leq x \leq 1$.


Analogo a lo hecho en el lenguaje fuente, los tipos de expresión no tipada y
tipada se definen parametrizando el arbol polimórfico con los pares de tipos
\verb|(unit,unit)| y \verb|(simple,distr)|, respectivamente.

\begin{ocamlblock}[caption={AST destino (extracto)}, label={lst:trgt-ast}]
  type ('s,'d) t =
  ...
| Sascr        of fsimple * ('s,'d) t * fsimple * 's annot
| Dascr        of fdistr  * ('s,'d) t * fdistr  * 'd annot
| Choice       of ('s,'d) t * ('s,'d) t * fprob * fprob * Form.t * 'd annot
| Let          of string  * ('s,'d) t * ('s,'d) t list * 's list * 'd annot

type untyped_term = (unit,unit) t
type typed_term   = (fsimple,fdistr) t
\end{ocamlblock}

\subsubsection{Valores}

En el módulo \verb|Target.Val| se definen los tipos de datos para los
valores retornados por el intérprete. El tipo de valor más básico es
el valor crudo, que puede ser un número, un booleano o una clausura (el
valor de una expresión lambda). Este último valor conserva en su
estructura el ambiente de definición, por lo que las funciones en esta
implementación tiene \emph{scope léxico}\footnote{definir}.

El tipo \verb|value| corresponde al valor que figura asociado a una
probabilidad en las distribuciones de valor retornadas por el
intérprete. El constructor de valor \verb|Ascr| representa una
ascripción de tipo simple (tercer argumento) sobre el valor crudo
(segundo argumento) contenido en el valor. El primer argumento es otro tipo
simple que funciona como evidencia de la ascripción.

La distribucion de valores es definida mediante el mismo tipo
paramétrico \verb|Formula.Distribution.dis| utilizado para definir los
tipos de fórmula de distribución. En este caso el tipo \verb|dis|
recibe como parámetro el tipo \verb|value|, lo que resulta en una
distribución simbólica sobre valores del lenguaje, con una fórmula que
restringe sus variables.

\begin{ocamlblock}[caption={Valor Crudo, Valor, y Distribucion de Valores}, label={lst:val}]
  (* Raw Value *)
  type raw_value = 
  | Num of float * val_id 
  | Bool of bool * val_id
  | Closure of (string * simple * typed_term * env) * val_id

  (* Value *)
  and value = 
  | Ascr of simple * raw_value * simple 
  | Error of simple
  
  (* Distribution of Values *)
  type distr_value = value Formula.Distribution.dis
\end{ocamlblock}

\subsection{Rutinas}

\subsubsection{Lexing y Parsing}

El análisis léxico del intérprete fue implementado haciendo uso de la
librería generadora de analizadores léxicos Ocamllex. Los
archivos implementados como artefactos para que dicha librería genere
el módulo \verb|Lexer| son:
\begin{itemize}
\item Un archivo con la implementación de un tipo de dato algebraíco
  denominado \verb|token|, cuyos constructores son una representación
  abstracta de las unidades léxicas que se espera encontrar en el
  texto. La mayoría de los constructores no contiene más que dos
  argumentos de información posicional: el inicio y el final del
  lexema en el texto. En cambio constructores como \verb|NUM| (número) o
  \verb|IDENT| (identificador) portan la información específica que le da
  a un token su expresividad. El primer constructor porta un número de punto
  flotante, y el segundo un string (el identificador o variable)
\item Un archivo con la definición de una regla que retorna un token,
  mediante \emph{pattern matching} sobre expresiones regulares, que se
  comparan con los lexemas encontrados en el texto. En casos como el de
  reconocimiento de un número en el texto, el número se pasa como un argumento
  al constructor del token \verb|NUM|, además de la información posicional.
  El caso del token \verb|LAMBDA| muestra que más de un lexema puede ser asociado
  un token.
\end{itemize}

\begin{ocamlblock}[caption={Implementación de \texttt{token} y reconocimiento de lexemas},label={lst:imp-parse}]
  type token =
  ...
  | NUM        of (float  * Lexing.position * Lexing.position)
  | IDENT      of (string * Lexing.position * Lexing.position)
  | LAMBDA     of (Lexing.position * Lexing.position)
  | TRUE       of (Lexing.position * Lexing.position)

  rule token = parse
  ...
  | decimal {
    let sp = Lexing.lexeme_start_p lexbuf in
    let ep = Lexing.lexeme_end_p   lexbuf in
    NUM (Float.of_string decimal, sp, ep)
  }
  ...
  | "lambda" | "lambda!" { 
    let sp = Lexing.lexeme_start_p lexbuf in
    let ep = Lexing.lexeme_end_p   lexbuf in
    LAMBDA(sp, ep)
  }
\end{ocamlblock}

El parser se implementó con la librería Menhir, a partir de un archivo
\verb|parser.mly| donde se declara la gramática del lenguaje. Esta
gramática describe cómo una secuencia de tokens (producidos por el
lexer) se transforma en expresiones del lenguaje: cada token incluye sus
posiciones (de tipo \verb|Lexing.position|), de modo que el parser puede
construir recursivamente el AST fuente con los constructores provistos
por \verb|Source.Term| y, al mismo tiempo, registrar en cada nodo su
rango posicional dentro del código fuente.  Durante el parsing, algunas
reglas de producción, como la de probabilidad gradual y la de tipos de
distribución, se someten a chequeos semánticos para verificiar su
validez.

\subsubsection{Chequeo Estático Fuente}
\TODO{Estoy revisando mi borrador para esta implementación}
\subsubsection{Operaciones de Tipo}

Las operaciones de tipo más relevantes que se implementan en este
intérprete son las que determinan la consistencia, precisión, y el
ínfimo entre dos tipos. Las siguientes razones justifican limitar la
exposición a las operaciones entre tipos de fórmula:
\begin{itemize}
\item La única operación que se usa en tipos graduales es la consistencia
\item Las tipos graduales de distribución se deben convertir antes a tipos de fórmula
  para determinar sus consistencia, por lo que la operación de consistencia para tipos
  graduales de distribución no es más que convertirlos a tipos de fórmula y luego aplicar
  la versión de fórmula de la operación entre los tipos convertidos
\item La consistencia entre tipos graduales simples se implementa de manera análoga
  entre tipos simples de fórmula.
\end{itemize}

Para tipos simples o de distribución, las operaciones de consistencia
y precisión toman dos tipos de la misma categoría
(\verb|(simple,simple)| o \verb|(distr,distr)|) y retornan un valor de
tipo \verb|bool|, que indica el valor lógico de la propocisión. En el
caso del meet, se retorna un tipo \verb|'t option|, donde \verb|'t| es
el tipo de los argumentos: Por ejemplo si los argumentos son de tipo
\verb|simple|, entonces el resultado puede ser \verb|Some Real|, en el
caso de que el resultado exista y sea \verb|Real|, y \verb|None| en
caso contrario.  Las versiones de estas operaciones para tipos
simples son una traducción directa de la formulación original del
lenguaje, codificando en casos de pattern matching los resultados
para cada par de tipos.

La consistencia simple recibe dos tipos simples \verb|t1| y
\verb|t2|. El pattern matching sobre el par \verb|(t1,t2)| retorna
\verb|true| para los pares \verb|(Any,_)| y \verb|(_,Any)|, donde
\verb|_| es cualquier tipo. También se retorna \verb|true| para los
casos \verb|(Real,Real)| y \verb|(Bool,Bool)|. El caso del par
\verb|(Fun(s1,d1),Fun(s2,d2))| retorna el valor \verb|b1 && b2|,
donde \verb|b1| es la consistencia simple entre \verb|s1| y \verb|s2|,
y \verb|b2| es la consistencia de distribución entre \verb|d1| y
\verb|d2|.

La operación de precisión retorna \verb|true| para los pares
\verb|(Real,Real)| y \verb|(Bool,Bool)|. El par \verb|(_,Any)| retorna
\verb|true| (todo tipo es tanto o más preciso que true). Siguiendo el
patrón de la consistencia, la precisión retorna
\verb|(Fun(s1,d1),Fun(s2,d2))| si la precisión simple entre s1 y s2, y
la precisión de distribución entre d1 y d2 retornan \verb|true|.


La operación de meet simple retorna un tipo opcional: si entre los tipos argumento está definida la relación de precisión, se retorna el constructor \verb|Some| con el tipo más preciso. Si los dos tipos no son comparables, la operación no está definida y se retorna \verb|None|. Concretamente si uno de los dos tipos es el desconocido \verb|Any|, se retorna por defecto \verb|Some| con el otro tipo. Para los tipos \verb|(Real,Real)| y \verb|(Bool,Bool)| se retorna \verb|Real| y \verb|Bool|, respectivamente.
Siguiendo el patrón recursivo de las operaciones anteriores,
para el par de tipos \verb|(Fun(s1,d1),Fun(s2,d2))| se retorna \verb|Some Fun(s3,d3)|, donde \verb|s3| es el tipo obtenido del meet simple entre \verb|s1| y \verb|s2|, y \verb|d3| es el tipo en el meet de distribución entre \verb|d1| y \verb|d2|. Si cualquiera de los dos meet recursivos retorna \verb|None|, entonces el meet entre funciones también lo hace.

Para los tipos de distribución de fórmula, las tres operaciones (\verb|consistency|, \verb|precision| y \verb|meet|) siguen un mismo patrón, basado en construir un acoplamiento simbólico entre las probabilidades de ambas distribuciones. Si $d_1=(\{(s_i,p_i)\}_i,\Phi_1)$ y $d_2=(\{(t_j,q_j)\}_j,\Phi_2)$, se crea una matriz de variables frescas $\alpha_{ij}$ (módulo \verb|Formula.Coupling|) y se impone: (i) \emph{compatibilidad por celda} ($\alpha_{ij}$ debe ser nula cuando $s_i$ no es compatible con $t_j$), (ii) \emph{marginales} ($p_i=\sum_j \alpha_{ij}$ y $q_j=\sum_i \alpha_{ij}$), y (iii) las fórmulas originales $\Phi_1$ y $\Phi_2$. El solver decide la satisfacibilidad del sistema; en consistencia/precisión eso determina el valor booleano, y en meet, si es satisfacible, el soporte del resultado se obtiene filtrando las celdas definidas (las que retornan \verb|Some|) y asociándoles su peso $\alpha_{ij}$.

Para los tipos de distribución de fórmula, las tres operaciones siguen
un mismo patrón de diseño e implementación (ilustrado en el algoritmo \ref{alg:typeop}), basado en la construcción
de un coupling $\coupling$, es decir una matriz de variables. La
primera tarea es determinar qué variables $\alpha_{ij}$ deben ser
iguales a 0 por la incompatibilidad de los tipos simples $\sigma_i$ y
$\delta_j$ de las distribuciones en cuestion. Para discernir esto se
define, a modo de \emph{máscara}, la matrix $R$ con entrada $i,j$ el
resultado \verb|meet_simple|$(\sigma_i,\delta_j)$ en el caso de la
operación meet, y el valor lógico \verb|rel|$(\sigma_i,\delta_j)$
cuando \verb|rel| es la consistencia o la precisión. De esta manera,
las funciones \verb|Coupling.alpha_option_formula| y
\verb|Coupling.alpha_bool_formula|, toman como argumento el par
$(R,\coupling)$ y retornan la fórmula que contiene la condición
correcta para cada $\alpha_{ij}$: Si
\verb|meet_simple|$(\sigma_i,\delta_j)$ está definido, o
\verb|rel|$(\sigma_i,\delta_j)$ es verdadero, entonces
$0<=\alpha_{ij}<= 1$ (es una variable libre) dentro de la fórmula. De
lo contrario, $\alpha_{ij} = 0$. La segunda tarea es imponer a cada
fila de la matriz $\coupling$ sumar lo mismo que su probabilidad
simbólica correspondiente, es decir $p_i = \sum_j\alpha_{ij}$. También
se debe imponer una condición análoga para las columnas y las
probabilidades $q_j$. Las formulas que codifican estas restricciones
para filas y columnas de $\coupling$ son retornadas
\verb|Coupling.row_sum_formula| y \verb|Coupling.col_sum_formula|
respectivamente, que toman como argumento el arrego de probabilidades
correspondiente y $\coupling$. La fórmula $\Phi$ que agrupa en
conjunción la fórmula de restricciones individuales sobre variables
$\alpha_{ij}$, las fórmulas de suma de filas y columnas, y las dos
fórmulas $\Phi_1$ y $\Phi_2$, heradadas de las distribuciones $\mu$ y
$\nu$, es la fórmula cuya solubilidad indica si el meet está definido,
o si las relaciones son verdaderas. La función
\verb|Solver.Z3solver.solve| es aplicada sobre la fórmula $\Phi$, y si
la fórmula tiene solución, entonces en este punto la consistencia y
precisión retornan \verb|true|. De lo contrario retornan
\verb|false|. Si $\Phi$ tiene solución y se trata de meet, entonces la
función \verb|Coupling.filter_alpha_option| retorna una distribución
con las entradas no nulas en la matriz $R$ asociadas a sus pesos
correspondientes en la matriz $\coupling$. Esta distribución y la
fórmula $\Phi$ son envueltos en \verb|Some| y retornados. En cambio si
la fórmula $\Phi$ no tiene solución, se retorna \verb|None|


\begin{algorithm}[H]
\caption{Patrón común para operaciones sobre tipos de distribución}
\label{alg:typeop}
\KwIn{
  $\mu=([(\sigma_i,p_i)]_i,\Phi_1)$, $\nu=([(\delta_j,q_j)]_j,\Phi_2)$,
  $\mathtt{op}\in\{\mathtt{Cons},\mathtt{Prec},\mathtt{Meet}\}$
}
\KwOut{\texttt{bool} (Cons/Prec) o \texttt{Option}$(D,\Phi)$ (Meet)}

$\coupling \leftarrow \texttt{Coupling.create}(m,n)$ \tcp*{Coupling con variables frescas $\alpha_{ij}$}

\lIf{$\mathtt{op}=\mathtt{Meet}$}
{$R \leftarrow [\,\texttt{meet\_simple}(\sigma_i,\delta_j)\,]_{ij}$}

\lElse{$R \leftarrow [\,\texttt{rel}(\sigma_i,\delta_j)\,]_{ij}$ \tcp*{\texttt{rel} es consistencia o precisión}}

\BlankLine


$\Phi \leftarrow \Phi_1 \land \Phi_2$ \tcp*{$\Phi$ es la fórmula a resolver}
\lIf{$\mathtt{op}=\mathtt{Meet}$}{$\Phi \leftarrow \Phi \land \texttt{Coupling.alpha\_option\_formula}(\coupling,R)$}

\lElse{$\Phi \leftarrow \Phi \land \texttt{Coupling.alpha\_bool\_formula}(\coupling,R)$}
$\Phi \leftarrow \Phi \land \texttt{Coupling.row\_sum\_formula}([p_i]_i,\coupling)$\;
$\Phi \leftarrow \Phi \land \texttt{Coupling.col\_sum\_formula}([q_j]_j,\coupling)$\;
\BlankLine

\If{\texttt{solve}$(\Phi)$}{
  \uIf{$\mathtt{op}\neq\mathtt{Meet}$}{
    \Return \texttt{true}\;
  }
  \Else{
    $D \leftarrow \texttt{Coupling.filter\_alpha\_option}(R,\coupling)$;
    
    $\xi \leftarrow (D,\Phi)$;
    
    \Return \texttt{Some} $\xi$\;
  }
}
\Else{
  \uIf{$\mathtt{op}\neq\mathtt{Meet}$}{
    \Return \texttt{false}\;
  }
  \Else{
    \Return \texttt{None}\;
  }
}
\end{algorithm}


\subsubsection{Resolución de Fórmulas}
\label{sec:impl-rutinas-solver}

La función \verb|Solver.Z3solver.solve| (algoritmo \ref{alg:solve}) es el punto de conexión del
intérprete con el solucionador Z3. Esta función recibe una lista de
fórmulas (de tipo \verb|Formula.Form.t|) y retorna un valor lógico:
\verb|true| si el sistema tiene una solución en los números reales, y
\verb|false| si no existe una solución.

En la implementación de \verb|solve|, la lista de fórmulas recibida
como argumento se transforma a una lista de expresiones nativas de la
librería \verb|Z3|. Esta transformación se realiza mediante la función
auxiliar \verb|formula_to_z3expr| (algoritmo \ref{alg:formula-to-z3expr}), que traduce cada fórmula del
intérprete a una expresión de Z3.

Un aspecto importante de esta traducción es el \emph{contexto} de Z3.
La función \verb|formula_to_z3expr| toma como argumento una fórmula y un
contexto \verb|ctx|. Este contexto agrupa la información necesaria para
construir y tipar expresiones dentro de una invocación al solucionador,
por lo tanto todas las transformaciones de fórmulas dentro de
\verb|solve| deben usar un mismo contexto.

En particular, \verb|formula_to_z3expr| primero utiliza
\verb|Algebra.of_formula| para obtener una restricción expresada en un tipo de datos interno del módulo \verb|Solver| a partir
de la fórmula, y luego utiliza \verb|Z3expr.of_constr| para convertir
dicha restricción en una expresión nativa de \verb|Z3|. El resultado
corresponde a igualdades/desigualdades cuyas variables son heredadas de
las probabilidades simbólicas presentes en la fórmula original.



\begin{algorithm}[H]
\caption{\texttt{formula\_to\_z3expr}}
\label{alg:formula-to-z3expr}
\KwIn{$\Phi$ : \texttt{Formula.Form.t}, $ctx$ : \texttt{Z3.context}}
\KwOut{$e$ : \texttt{Z3expr.t}}

$c \leftarrow \texttt{Algebra.of\_formula}(\Phi)$;

$e \leftarrow \texttt{Z3expr.of\_constr}(c,ctx)$ \tcp*{genera una expresión nativa de Z3}
\Return $e$\;
\end{algorithm}


Una vez obtenida la lista de expresiones de Z3, \verb|solve| intenta
resolver el sistema con dos solucionadores. El primero es un
solucionador de Z3 configurado para la lógica
\texttt{QF\_LRA}\footnote{\textbf{QF} significa \emph{quantifier-free}
  (sin cuantificadores). \textbf{LRA} es \emph{linear real
    arithmetic}: aritmética real lineal, es decir, restricciones sobre
  reales sin no-linealidades (p.\ ej., no hay productos entre
  variables como $x\cdot y$).}. Si el chequeo retorna \verb|SAT|
(\emph{satisfiable}), entonces \verb|solve| retorna \verb|true|,
indicando que el sistema tiene solución. Por otro lado, si el chequeo
retorna \verb|UNSAT| (\emph{unsatisfiable}), entonces \verb|solve|
retorna \verb|false|.

Si el chequeo resulta en \verb|UNKNOWN|, entonces se prueba con un
segundo solucionador de Z3 configurado para la lógica
\texttt{QF\_NRA}\footnote{\textbf{NRA} es \emph{nonlinear real
    arithmetic}: aritmética real no-lineal, es decir, permite
  no-linealidades (p.\ ej., productos entre variables como $x\cdot y$
  o potencias como $x^2$).}.  En este caso, \verb|solve| retorna
\verb|true| si el chequeo encuentra una solución, y retorna
\verb|false| en cualquier otro caso, para señalar que no existe una
solución en los reales para el sistema.


El resultado de \verb|solve| es recibido en un contexto donde se decide
qué hacer con su valor retornado pero, típicamente, \verb|true|
significa que una relación entre distribuciones se verifica o que una
operación entre distribuciones está bien definida. Por otro lado,
\verb|false| suele significar que se levantará una excepción.

\begin{algorithm}[H]
\caption{\texttt{solve}}
\label{alg:solve}
\KwIn{$[\Phi_i]_i$ : lista de \texttt{Formula.Form.t}}
\KwOut{\texttt{bool}}

$\mathit{cfg} \leftarrow \{(\texttt{"model"},\texttt{"true"})\}$\tcp*{configuración}
$ctx \leftarrow \texttt{Z3.mk\_context}(\mathit{cfg})$\tcp*{contexto}

$E \leftarrow [\,\texttt{formula\_to\_z3expr}(\Phi_i,ctx)\,]_i$\tcp*{Lista de expresiones}

$s_1 \leftarrow \texttt{mk\_solver}(ctx,\texttt{"QF\_LRA"})$\tcp*{Solucionador lineal}
\texttt{add}$(s_1, E)$;

$r_1 \leftarrow \texttt{check}(s_1)$\;

\If{$r_1 = \texttt{SAT}$}{\tcp*{$\Phi$ tiene solución}
  \Return \texttt{true}\;
}
\ElseIf{$r_1 = \texttt{UNSAT}$}{\tcp*{$\Phi$ NO tiene solución}
  \Return \texttt{false}\;
}
\Else{ \tcp*{$r_1 = \texttt{UNKNOWN}$}
  $s_2 \leftarrow \texttt{mk\_solver}(ctx,\texttt{"QF\_NRA"})$\tcp*{Solucionador NO lineal}
  \texttt{add}$(s_2, E)$\;
  $r_2 \leftarrow \texttt{check}(s_2)$\;

  \If{$r_2 = \texttt{SAT}$}{
    \Return \texttt{true}\;
  }
  \Else{
    \Return \texttt{false}\;
  }
}
\end{algorithm}




\subsubsection{Elaboración}
\label{sec:impl-rutinas-elaboracion}

La elaboración corresponde a la traducción desde el AST fuente
tipado (producido por \verb|Typing.Typecheck.typecheck|) hacia
un AST destino no tipado, implementada en el módulo
\verb|Target.Elaborate| mediante la rutina recursiva
\verb|Target.Elaborate.elaborate|. El objetivo de esta fase no es
volver a inferir tipos, sino \emph{reformular} el programa en un
lenguaje destino que haga explícitos los chequeos dinámicos y la
infraestructura simbólica (probabilidades simbólicas y fórmulas)
requerida por el tiempo de ejecución.

En términos de implementación, la elaboración tiene efectúa los
siguientes cambios: \textbf{1.} Cada nodo destino (incluyendo nodos
nuevos introducidos por la transformación) hereda los metadatos del
nodo fuente correspondiente, de modo que los errores reportados en
fases posteriores puedan localizarse con precisión.  \textbf{2.} Los
tipos graduales anotados en el AST fuente se transforman a tipos de
fórmula del lenguaje destino (mediante las fórmulas
\lstinline|lift_simple| y \lstinline|lift_distr|), lo cual hace
explícito el componente simbólico.  \textbf{3.} Cada ascripción del
lenguaje destino (\lstinline|Sascr| / \lstinline|Dascr|) porta una
evidencia, computada como el ínfimo entre el tipo inferido de la
subexpresión y el tipo esperado.

Un aspecto importante es que, dado que el AST fuente ya fue chequeado
estáticamente, la elaboración asume que los \emph{meet} necesarios
están definidos. En consecuencia, un fallo al calcular evidencia se
interpreta como un error interno de la implementación (o una violación
de invariantes), y no como un error atribuible al programa del
usuario.

El Algoritmo~\ref{alg:elab} muestra un extracto de casos
representativos.  Los casos omitidos (por ejemplo, operadores,
\lstinline|let|, ascripción de distribución, etc.)  siguen el mismo
patrón: recursión, levantamiento de tipos, computo de evidencia
mediante {meet}, y reconstrucción del nodo destino preservando
metadatos.


\begin{algorithm}[H]
\caption{Elaboración: \texttt{Target.Elaborate.elaborate} (extracto)}
\label{alg:elab}
\KwIn{$t$ (AST fuente tipado)}
\KwOut{$u$ (AST destino no tipado)}
$md \leftarrow \mathtt{meta}(t)$; \tcp*{metadatos del nodo fuente}
$ann \leftarrow \mathtt{Annot}((),md)$; \tcp*{anotación destino (sin tipo, solo metadatos)}

\If{$t = \mathtt{Num}(n,\_)$}{
  \Return $\mathtt{Sascr}(Real,\ \mathtt{Num}(n,ann),\ Real,\ ann)$\;
}

\ElseIf{$t = \mathtt{Var}(x,\_)$}{
  \Return $\mathtt{Var}(x,ann)$\;
}



\ElseIf{$t = \mathtt{Dascr}(m,\nu,\_)$}{
  $\mu \leftarrow \texttt{distrType}(m)$\tcp*{tipo gradual de expresión fuente}
  $\mu' \leftarrow \mathtt{lift\_distr}(\mu)$; $\ \nu' \leftarrow \mathtt{lift\_distr}(\nu)$\tcp*{tipos grad -> tipos de fórm.}
  $m' \leftarrow \mathtt{elaborate}(m)$\tcp*{elaboración de m}
  $\epsilon \leftarrow \mathtt{meet\_distr}(\mu',\nu')$\tcp*{se asume definido}
  \Return $\mathtt{Dascr}(\epsilon,\ m',\ \nu',\ ann)$\;
}

\ElseIf{$t = \mathtt{Choice}(p,m,n,\_)$}{
  $\mu \leftarrow \mathtt{distrType}(m)$; $\ \nu \leftarrow distrType\mathtt{}(n)$;
  
  $m' \leftarrow \mathtt{elaborate}(m)$; $\ n' \leftarrow \mathtt{elaborate}(n)$;
  
  $\mu' \leftarrow \mathtt{lift\_distr}(\mu)$; $\ \nu' \leftarrow \mathtt{lift\_distr}(\nu)$\;

  $(\alpha,\Phi_\alpha) \leftarrow \mathtt{lift\_gprob}(p)$\tcp*{prob. simbólica + restricciones}
  $(\omega,\Phi_{\omega}) \leftarrow \mathtt{complement}(\alpha)$\;
  $\Phi \leftarrow \Phi_\alpha \land \Phi_{\omega} \land (\alpha + \omega = 1)$;

  $\tau_1 \leftarrow \mathtt{add\_formula}(\alpha \cdot \mu' + \omega\cdot\nu', \Phi )$\;

  $\tau_2 \leftarrow \mathtt{lift\_distr}(\mathtt{distrType}(t))$\tcp*{tipo gradual inferido del nodo}
  $\epsilon \leftarrow \mathtt{meet\_distr}(\tau_1,\tau_2)$\tcp*{se asume definido}

  $c \leftarrow \mathtt{Choice}(m',n',(\alpha,\Phi_\alpha),(\omega,\Phi_\omega),\Phi,ann)$\;
  \Return $\mathtt{Dascr}(\epsilon,\ c,\ \tau_2,\ ann)$\;
}

\tcp*{otros casos}

\Else{

   \textbf{Error} \;
}
\end{algorithm}


\subsubsection{Chequeo Estático Destino}
\label{sec:impl-rutinas-tc-destino}

\TODO{También tengo un borrador en revisión, casi listo}


\subsubsection{Evaluación}
